\documentclass{article}
\usepackage[utf8]{inputenc}
\usepackage{amsmath,amssymb}
\usepackage[most]{tcolorbox}
\usepackage{geometry}
\geometry{margin=2.5cm}

\newcommand{\step}[1]{\par\noindent\hangindent=3.5em\hangafter=1%
  \makebox[3.5em][l]{\textbf{#1.}}}
\newcommand{\steptag}[1]{\hfill{\small\textsf{[#1]}}}

\newtcolorbox{proofbox}[1]{
  colback=gray!5, colframe=gray!60,
  fonttitle=\bfseries, title={#1},
  breakable, enhanced,
  left=4pt, right=4pt, top=4pt, bottom=4pt,
  before skip=10pt, after skip=10pt
}

\begin{document}

\begin{proofbox}{Group-input polar-domain membership: there exist universa...}
\step{1} Group-input polar-domain membership: there exist universal C\_grp, C\_pol >= 1, kappa\_pol in (0, 1/2] such that for every finite-dimensional exact-unit epsilon\_r-C*-algebra and delta > 0 satisfying C\_pol*(epsilon\_r + delta) <= kappa\_pol and C\_grp*epsilon\_r < delta - C\_pol*(epsilon\_r*delta + delta\textasciicircum{}2), the inverse u\_delta:S\_delta -> calU of the polar map is defined at U bold-dot V and U\textasciicircum{}dagger for every U, V in calU; moreover, U bold-dot V and U\textasciicircum{}dagger each have a right inverse. \steptag{validated} \\[6pt]
\step{1.1} Invoke the registered external lem-stage1-polar-retraction and choose its universal constants C\_pol >= 1 and kappa\_pol in (0,1/2]; set the further universal witness C\_grp:=4. Fix an arbitrary finite-dimensional exact-unit epsilon\_r-C*-algebra, delta>0 satisfying the two root guards, and U,V in calU; abbreviate e:=epsilon\_r and t:=delta-C\_pol*(e*delta+delta\textasciicircum{}2). The external inner inclusion calU\_t subseteq S\_delta reduces the root to proving U bold-dot V and U\textasciicircum{}dagger belong to calU\_t, including the right-inverse clause in the definition of calU\_t. \steptag{validated} \\[4pt]
\step{1.2} Guard arithmetic in the setup of node 1.1: C\_pol*(e+delta)<=kappa\_pol<=1/2 and C\_pol>=1 give s:=e+delta<=1/2 and delta<=1/2; the second guard with C\_grp=4 gives 4e<t, hence t>0. Also t=delta-C\_pol*delta*s<=delta*(1-s). Since e<t<=delta*(1-e-delta), write delta=s-e. The function (s-e)*(1-s) is increasing for e<=s<=1/2 because its derivative is 1+e-2s>=e>=0; therefore e<(1/2-e)/2=1/4-e/2, so e<1/6. These are the only smallness consequences used below. \steptag{validated} \\[4pt]
\step{1.3} Common multiplier estimate, using node 1.2 explicitly. For every W in calU, def-approximate-unitary-space gives W\textasciicircum{}dagger bold-dot W=J and a right inverse, while def-epsilon-cstar-algebra gives (1-e)||W||\textasciicircum{}2<=1 and ||L\_\{W\textasciicircum{}dagger\}L\_W-I||<=e||W||\textasciicircum{}2<=e/(1-e)<1/5. Thus L\_W is injective and, because calX is finite-dimensional, bijective. Moreover ||L\_\{W\textasciicircum{}dagger\}||<=(1+e)||W||<=7/5 and ||L\_\{W\textasciicircum{}dagger\}L\_W Z||>=(4/5)||Z||, so ||L\_W Z||>=(4/7)||Z|| and ||L\_W\textasciicircum{}\{-1\}||<=7/4. All inequalities use only e<1/6 from node 1.2 and exact unitality. \steptag{validated} \\[4pt]
\step{1.4} Product defect, using nodes 1.2 and 1.3. Put X:=U bold-dot V. The star rule gives X\textasciicircum{}dagger=V\textasciicircum{}dagger bold-dot U\textasciicircum{}dagger. Insert ((V\textasciicircum{}dagger bold-dot U\textasciicircum{}dagger) bold-dot U) bold-dot V and (V\textasciicircum{}dagger bold-dot (U\textasciicircum{}dagger bold-dot U)) bold-dot V between X\textasciicircum{}dagger bold-dot X and J. The two approximate-associativity errors are each at most e*(1+e)||U||\textasciicircum{}2||V||\textasciicircum{}2, and the final expression is (V\textasciicircum{}dagger bold-dot J) bold-dot V=V\textasciicircum{}dagger bold-dot V=J. Hence ||X\textasciicircum{}dagger bold-dot X-J||<=2e(1+e)/(1-e)\textasciicircum{}2<=4e<t<2t, where e<1/6 and t>4e are exactly node 1.2. \steptag{validated} \\[4pt]
\step{1.5} Product right inverse, using nodes 1.2 and 1.3. Approximate associativity gives ||L\_\{U bold-dot V\}-L\_U L\_V||<=e||U||||V||<=e/(1-e)<1/5. The operator A:=L\_U L\_V is invertible and ||A\textasciicircum{}\{-1\}||<=||L\_V\textasciicircum{}\{-1\}||||L\_U\textasciicircum{}\{-1\}||<=49/16, so ||A\textasciicircum{}\{-1\}(L\_\{U bold-dot V\}-A)||<49/80<1. The Neumann lemma makes L\_\{U bold-dot V\}=A[I+A\textasciicircum{}\{-1\}(L\_\{U bold-dot V\}-A)] invertible. Therefore R\_X:=L\_\{U bold-dot V\}\textasciicircum{}\{-1\}J satisfies (U bold-dot V) bold-dot R\_X=J and is a right inverse of U bold-dot V. \steptag{validated} \\[4pt]
\step{1.6} Adjoint input, using nodes 1.2 and 1.3. Def-approximate-unitary-space gives U\textasciicircum{}dagger bold-dot U=J, so U is a right inverse of U\textasciicircum{}dagger. Let R be the right inverse of U supplied by that same definition. Since L\_U is bijective by node 1.3, R=L\_U\textasciicircum{}\{-1\}J and ||R||<=7/4. Exact unitality and one associator estimate give ||U\textasciicircum{}dagger-R||=||U\textasciicircum{}dagger bold-dot(U bold-dot R)-(U\textasciicircum{}dagger bold-dot U) bold-dot R||<=e||U||\textasciicircum{}2||R||. Consequently ||U bold-dot U\textasciicircum{}dagger-J||=||U bold-dot(U\textasciicircum{}dagger-R)||<=(1+e)e||U||\textasciicircum{}3||R||<=4e<t<2t; the numerical coefficient is at most (7/6)*(6/5)\textasciicircum{}3*(7/4)=441/125<4 by e<1/6. Thus U\textasciicircum{}dagger has the required strict defect bound and an explicit right inverse. \steptag{validated} \\[4pt]
\step{1.7} Membership conclusion. Nodes 1.4 and 1.5 give ||(U bold-dot V)\textasciicircum{}dagger bold-dot(U bold-dot V)-J||<2t together with a right inverse, so def-approximate-unitary-space gives U bold-dot V in calU\_t. Node 1.6 gives the same two defining properties for U\textasciicircum{}dagger, hence U\textasciicircum{}dagger is in calU\_t. This also records the right inverses required separately by the root contract. \steptag{validated} \\[4pt]
\step{1.8} Apply the registered external lem-stage1-polar-retraction with the constants fixed in node 1.1. Its guard is the first root guard and node 1.2 gives t>0; its inner inclusion calU\_t subseteq S\_delta and node 1.7 therefore put U bold-dot V and U\textasciicircum{}dagger in S\_delta. Hence the inverse component u\_delta:S\_delta->calU is defined at both points. Their right inverses were constructed in nodes 1.5 and 1.6. Since the algebra, delta, U, and V were arbitrary, the witnesses C\_grp=4 together with C\_pol,kappa\_pol prove the root contract. \steptag{validated} \\[4pt]
\step{1.9} Endpoint-safe replacement for the endpoint-sensitive portions of nodes 1.3--1.6 (those earlier strict chains are not used here), using only validated nodes 1.1 and 1.2 and the registered definitions. For W in calU, W\textasciicircum{}dagger bold-dot W=J and (1-e)||W||\textasciicircum{}2<=1. Exact unitality and the associator axiom give ||L\_\{W\textasciicircum{}dagger\}L\_W-I||<=e||W||\textasciicircum{}2<=e/(1-e)<1/5. Hence, weakening all consequent bounds non-strictly, ||L\_\{W\textasciicircum{}dagger\}L\_W Z||>=(4/5)||Z|| and ||L\_\{W\textasciicircum{}dagger\}||<=(1+e)||W||<=7/5, so ||L\_W Z||>=(4/7)||Z||; finite dimensionality makes L\_W bijective and ||L\_W\textasciicircum{}\{-1\}||<=7/4. For X:=U bold-dot V, two associator estimates and U\textasciicircum{}dagger bold-dot U=V\textasciicircum{}dagger bold-dot V=J yield ||X\textasciicircum{}dagger bold-dot X-J||<=2e(1+e)/(1-e)\textasciicircum{}2<=4e<t<2t. Also ||L\_X-L\_U L\_V||<=e||U||||V||<=e/(1-e)<=1/5, while ||(L\_U L\_V)\textasciicircum{}\{-1\}||<=49/16; thus the perturbation norm is <=49/80<1, L\_X is invertible, and L\_X\textasciicircum{}\{-1\}J is a right inverse of X. Finally let R be the right inverse of U supplied by calU. The preceding bijectivity gives R=L\_U\textasciicircum{}\{-1\}J and ||R||<=7/4. Since U bold-dot R=J and U\textasciicircum{}dagger bold-dot U=J, one associator estimate gives ||U\textasciicircum{}dagger-R||<=e||U||\textasciicircum{}2||R||, hence ||U bold-dot U\textasciicircum{}dagger-J||<=(1+e)e||U||\textasciicircum{}3||R||<=4e<t<2t; U is a right inverse of U\textasciicircum{}dagger. (The coefficient bounds use only e<1/6 from node 1.2: ||W||<=6/5, 2(1+e)/(1-e)\textasciicircum{}2<=4, and (1+e)(6/5)\textasciicircum{}3(7/4)<=441/125<4.) Thus X and U\textasciicircum{}dagger lie in calU\_t with explicit right inverses. By the inner inclusion calU\_t subseteq S\_delta from node 1.1, u\_delta is defined at both. Every defect/operator estimate up to its constant multiple is <=, including when e=0; the only variable-dependent strict passage is 4e=C\_grp*e<t (followed by t<2t because t>0). This independently proves the root without the challenged strict inferences. \steptag{validated} \\[4pt]
\end{proofbox}

\end{document}
